\documentclass[dvipdfmx]{jsarticle}
\usepackage{graphicx}
\usepackage{url}

\title{ソフトウェア設計及び実験\\第7回（UI3）}
\author{学生番号 \ 6125010806　氏名 \ 清水 \ 統貴}
\date{2026年6月日提出}

\begin{document}
\maketitle

\section{課題の目的}
%
% 1. 目的
% 2. 原理・理論
% 3. 方法
% 4. 結果
% 5. 考察
% 6. 参考文献
% 7. 付録
本課題の目的は以下のとおりである．\\
• Joy-Conを扱うプログラミングの習得.\\
• 大規模コーディングの基本を学習.\\

% \section{原理}
% %

\section{課題1}
%
\subsection{joycon\_axis.c}
%
このプログラムはJoy-Conの6軸センサーの情報を取得するものである. 
取得したジャイロセンサーの情報を用いて, 傾きを計算しており, センサー
から取得した情報は角度ではなく, 角速度なので, それを時間で積分すること
で角度を出している.\\
\subsection{joycon\_ir.c}
%
このプログラムはJoy-Conの底面についているIRモーションカメラを使って画像を
撮影するものである. 撮影した画像は\verb+irdata.png+というファイルとして保存
される.\\
\subsection{joycon\_led.c}
%
このプログラムはJoy-Conの側面及びホームボタンの周りについているLEDを光らせるものである. 
側面のLEDは点灯, 点滅の切り替えをして, ホームボタンは光量や光る時間などを調整して, ゆっくり
明るくなり, ゆっくり暗くなるような動作をする.\\
\subsection{joycon\_nfc.c}
%
このプログラムでは, Joy-Conのスティックに搭載されているNFCリーダーを起動して, NFCカードなどを
読み取るものとなっている. 読み取った情報は画面に表示されるようになっている.\\
\subsection{joycon\_rumble.c}
%
このプログラムはJoy-Conに搭載されたHD振動を使って, Joy-Conを振動させたり, Joy-Conから音楽を流したり
するものである. Xボタンを押している間は, 一定の強さで振動し, Yボタンを押している間は, 周波数や振幅
など, より詳細な設定をした振動をし, ZRボタンを押すと, 音楽が流れるようになっている.\\
\subsection{joycon\_state.c}
%
このプログラムはJoy-Conの基本的な動作を読み取るものである. ボタンが押されている時にどこのボタンが押されているか
表示され, 別のボタンが押されるようになったらボタンの状態が変わったことが表示される. また, スティックがどの位置を
向いているかも表示される.\\
\subsection{Joy-Conライブラリの基本的な使い方}
%
まずプログラムの先頭にJoy-Conライブラリをインクルードする. 次に\verb+joycon_open()+でJoy-Conとの接続を開始する. 
そして, mainループの中で任意の関数を呼び出し, 各動作の処理を行う. 最後に\verb+joycon_close()+を呼び出し, 接続を
終了する.\\
% \section{課題2}
% %
% \section{感想}
% %


\begin{thebibliography}{99}
%
% \bibitem{chatGPT} Open AI : ChatGPT \url{https://chatgpt.com/ja-JP} (2026年4月30日閲覧).
% %
% \bibitem{claude} Anthropic : Claude \url{https://claude.ai/new} (2026年5月5日閲覧).
% %
\bibitem{gemini} Google : Gemini \url{https://gemini.google.com/?hl=ja}
%
\bibitem{imo} Imo : 6軸センサーから3軸回転の傾き角度算出とドリフト補正方法 \url{https://garchiving.com/calculate-angle-of-3axis-rotation-in-6axis-sensor/}
%
\bibitem{qiita} awesomecity : OpenCV　「画像のグレースケール化」 \url{https://qiita.com/awesomecity1998/items/448a3f91a489c8959d20}
%
\bibitem{nsi} NSI : 2値化とは？ \url{https://nsi-freak.com/binarization/#:~:text=2%E5%80%A4%E5%8C%96%E3%81%A8%E3%81%AF%E3%80%81%E7%94%BB%E5%83%8F%E3%82%92%E7%99%BD%E3%81%A8%E9%BB%92,%E3%80%8C%E9%96%BE%E5%80%A4%E3%80%8D%E3%81%A8%E8%A8%80%E3%81%84%E3%81%BE%E3%81%99%E3%80%82&text=%E3%83%A2%E3%83%8E%E3%82%AF%E3%83%AD%E7%94%BB%E5%83%8F%E3%81%AF0%EF%BD%9E,%E3%81%A8%E5%AE%9A%E7%BE%A9%E3%81%95%E3%82%8C%E3%81%A6%E3%81%84%E3%81%BE%E3%81%99%E3%80%82}
%
\bibitem{kuroro} 黒光俊秀 : OpenCVで使われるcannyとは?利用法からCanny法の理論などを徹底解説 \url{https://kuroro.blog/python/wOt3yEohr7oQt1qzif71/}
%
\end{thebibliography}
\end{document}